\documentclass[11pt]{ctexart}

\usepackage[a4paper,margin=1in]{geometry}
\usepackage{amsmath,amssymb,amsthm,mathtools}
\usepackage{algorithm}
\usepackage{algpseudocode}
\usepackage{hyperref}
\usepackage{enumitem}
\usepackage{booktabs}
\usepackage{graphicx}
\usepackage{float}

% 用 xelatex 编译本文件（ctex 宏包需要 xelatex/lualatex）：
%   xelatex nextval_unified_zh.tex
%   xelatex nextval_unified_zh.tex   % 第二遍用于交叉引用

\hypersetup{colorlinks=true,linkcolor=blue,citecolor=blue,urlcolor=blue}

\title{用于环检测的单调单指针锚点前进算法：\\
精确闭式、可证支配性,以及后继求值模型下的最优性}
\author{zx l}
\date{2026 年 6 月}

\newtheorem{theorem}{定理}[section]
\newtheorem{lemma}[theorem]{引理}
\newtheorem{proposition}[theorem]{命题}
\newtheorem{corollary}[theorem]{推论}
\newtheorem{definition}[theorem]{定义}
\newtheorem{remark}[theorem]{注}

\newcommand{\nextop}{\texttt{next}}
\newcommand{\nullnode}{\texttt{null}}
\newcommand{\oldc}{\texttt{old}}
\newcommand{\newc}{\texttt{newc}}

\renewcommand{\proofname}{证明}

\begin{document}

\maketitle

\begin{abstract}
给定一个起始结点与一个后继预言机 $f$,检测序列 $x,f(x),f^2(x),\dots$ 最终进入的
环,并恢复其尾长 $\mu$ 与周期 $\lambda$,这是经典问题,Floyd 的龟兔算法与 Brent
的二的幂分块算法都能以 $O(1)$ 内存解决。本文针对同一问题提出一种\emph{单调单指针
锚点前进}算法。它维护一个当前锚点、一个候选下一锚点,以及一个安全的定位锚点;以两
阶段方式探测,并采用一条失败压缩(``nextval'')规则,绝不重新扫描那些\emph{返回锚
点}检查已经失败过的区间。指针从不重走任何结点。在后继求值代价模型下,我们对该方法
给出一个统一而自洽的论述:\textbf{(i)} 正确性证明(线性时间、常数空间、安全定位);
\textbf{(ii)} 求值次数的\emph{精确闭式} $A(\mu,\lambda)=2(\mu+\lambda)+R$,其开销
$R\in\{1,2^{r_{\min}}-1\}$,并在 $\mu,\lambda\le 2000$ 的全部 $4{\times}10^{6}$ 个
格点上逐位验证;\textbf{(iii)} 一个\emph{被证明}的逐实例支配定理及其闭式平局集——
$A\le B$ 且等号成立\emph{当且仅当}
$(\mu,\lambda)\in\{(0,1),(1,1),(1,2)\}$,而 $A\le F$ 仅在 $(1,1)$ 取等;
\textbf{(iv)} 一个匹配的 $\mu+\lambda$ 下界,证明三种方法均为常数因子最优,且 $A$ 在
每个输入上都不超过该下界的 $3$ 倍;\textbf{(v)} 一个\emph{平均情形}常数
$c_A=\mathbb E[A]/\mathbb E[\mu+\lambda]\approx 2.18033$,经 Mellin 分析证明它是 $n$
的\emph{对数周期}函数,相比之下 Brent 约为 $3.08$、Floyd 约为 $3.5$;
\textbf{(vi)} 一个多锚点推广,刻画检测延迟与总工作量之间的 Pareto 前沿;以及
\textbf{(vii)} wall-clock 与 Pollard-$\rho$ 基准测试,定位出``更少的后继求值''何时
转化为真实加速的临界点。我们刻意把适用范围说清楚:渐近类 $\Theta(\mu+\lambda)$ 不变,
本方法比经典算法做更多的指针比较,其收益只有当单次后继求值昂贵时才兑现——也就是当年
促使 Brent 出现的 Pollard-$\rho$ 情形。
\end{abstract}

\section{引言}
\label{sec:intro}

许多计算都可归结为在三条约束下迭代一个确定映射:映射是\emph{黑盒}(只能观测到
$y=f(x)$,看不到其全局结构);只能做\emph{局部步进}(施加 $f$);只有\emph{有限内存}
(无法存下整条轨迹)。在这些约束下,有限状态空间上的轨迹
$x_0,x_1=f(x_0),\dots$ 最终是周期的:长度为 $\mu\ge 0$ 的尾巴接入长度为 $\lambda\ge 1$
的环(即 ``$\rho$'' 形)。恢复 $(\mu,\lambda)$ 是众多任务的核心原语:Pollard 的 $\rho$
因子分解与离散对数方法、伪随机数发生器的周期测试、模型检测中的非终止与套索检测、迭代
动力学中的极限环检测、以及路由环路或流去重检查。该问题也是经典的``在单链表中寻找环
入口'',其中 \nextop{} 指针即为后继预言机。在所有这些场景里,后继都是计量代价的单位;
而在其中若干场景(密码学映射、繁重的数值迭代)中,单次后继求值就主导了运行时间。

Floyd 的龟兔算法(经 Knuth 推广)与 Brent 的二的幂分块算法都以 $O(1)$ 内存求解,但每
单位 $\mu+\lambda$ 分别约花费 $3.5$ 与 $3.08$ 次后继求值。在局部预言机模型中,信息论
下界是 $\mu+\lambda$ 次求值:每个可达的不同结点都必须至少被触及一次
(定理~\ref{thm:lowerbound})。本文部分弥合了这一差距:所提方法的平均常数约为 $2.18$,
在\emph{每个}输入上都可证地支配两种经典算法,并且——不同于其前辈——它允许精确、可
验证的闭式分析。

该算法源于以下观察。在 Brent 风格的锚点方法中,一轮失败之后,下一轮的第一阶段可能恰好
重复失败轮第二阶段已扫描过的区间。若该重复区间针对相关锚点的\emph{返回}检查此前已经做
过,那么重新扫描它不可能产生新结果;于是我们\emph{压缩}该失败区间。方法维护三个特别的
指针:
\begin{itemize}[leftmargin=2em]
    \item $c$,当前锚点;
    \item $\newc$,由一轮第一阶段产生的候选下一锚点;
    \item $\oldc$,一个被证明位于环入口之前的定位锚点(或在边界情形 $\mu=0$ 时就是入口本身)。
\end{itemize}
定位锚点是第二个节省来源:因为 $\oldc$ 通常严格比表头更靠近入口,所以入口定位严格地比
Brent 更便宜。

设 $A(\mu,\lambda)$、$B(\mu,\lambda)$、$F(\mu,\lambda)$ 分别表示在尾长 $\mu$、周期
$\lambda$ 的环形输入上,所提算法、Brent 算法、Floyd 算法执行的后继求值次数。我们证明在
每个环形输入上 $A\le B$ 且 $A\le F$,给出 $A$ 的精确值,并精确刻画等号何时成立。

\paragraph{论断的适用范围。}
这并非主张普遍意义上的 wall-clock 优越。本方法会做额外的指针比较并维护额外状态;其优势
是在\emph{后继求值模型}下陈述的——其中每次操作 $p\leftarrow f(p)$ 计一单位代价,而比较、
赋值与整数算术免费。当后继求值占主导时(如函数迭代、Pollard-$\rho$ 因子分解或指针机
场景),这正是自然的模型。第~\ref{sec:wallclock}--\ref{sec:rho} 节在真实硬件上量化:一
次后继步必须昂贵到何种程度,模型层面的优势才会变成 wall-clock 上的优势。

\paragraph{贡献与组织。}
\begin{enumerate}[leftmargin=2em,itemsep=2pt]
    \item 一个常数空间、单调单指针的锚点前进算法(第~\ref{sec:algo} 节),配有一条探测
          调度不变量以证明跳过规则不会漏掉任何成功的返回锚点事件(第~\ref{sec:schedule}
          节),以及完整的正确性证明(第~\ref{sec:correctness} 节)。
    \item 求值次数的\emph{精确闭式}(定理~\ref{thm:closed-form}),在 $4{\times}10^{6}$
          个格点上验证(第~\ref{sec:closedform} 节)。
    \item 一个\emph{被证明}的、对 Brent 与 Floyd 的紧支配结果,带闭式平局集
          (第~\ref{sec:dominance} 节),把先前``在无穷多输入上严格''的陈述精化为一个
          有限、显式的等号集合。
    \item 一个匹配的 $\mu+\lambda$ 下界与常数因子最优性(第~\ref{sec:lowerbound} 节)。
    \item 平均情形分析:经 Mellin 残差得到 $c_A\approx 2.18033$ 作为 $n$ 的对数周期
          函数,并以蒙特卡洛验证(第~\ref{sec:average} 节)。
    \item 一个多锚点的空间--时间权衡,刻画检测/总工作量的 Pareto 前沿
          (第~\ref{sec:tradeoff} 节)。
    \item 可复现的操作计数、wall-clock 与 Pollard-$\rho$ 评测,定位廉价/昂贵后继的临界点
          (第~\ref{sec:empirical}--\ref{sec:rho} 节)。
\end{enumerate}

\section{模型与记号}
\label{sec:model}

设有限集 $S$($|S|=n$)上的映射 $f:S\to S$,并设从起点 $x_0$ 沿 $f$(等价地,沿
\nextop{} 指针)生成的序列为 $v_0,v_1,v_2,\dots$。存在唯一整数 $\mu\ge 0$(尾长)与
$\lambda\ge 1$(周期),使得
\[
v_{\mu+k+\lambda}=v_{\mu+k}\qquad\text{对一切 }k\ge 0,
\]
且对一切相异的 $i,j<\mu+\lambda$ 有 $v_i\ne v_j$。我们把每个绝对位置 $x\in\mathbb N$ 与
如下结点对应:
\[
N(x)=\begin{cases}x,&x<\mu,\\ \mu+((x-\mu)\bmod \lambda),&x\ge \mu,\end{cases}
\]
于是位置 $x,y$ 表示同一结点当且仅当 $N(x)=N(y)$。记 $k=\mu+\lambda$。

\begin{definition}[代价模型]
\label{def:cost}
一次计算的代价是后继求值 $p\leftarrow f(p)$ 的次数。指针比较、赋值、布尔测试与整数算术
不计入。支配性陈述仅针对环形输入作出;对无环输入,我们仅证明线性时间终止。
\end{definition}

\begin{remark}[绝对位置与指针值]
分析使用一条提升后的执行轨迹,其中每次后继求值都把移动指针推进一个绝对位置。在环形结构
中,不同的绝对位置可能表示同一具体结点,因此存储的指针值并不能唯一确定一个位置,但按时间
顺序的轨迹可以。下文所有关于区间的陈述都指这条提升轨迹。
\end{remark}

\section{算法}
\label{sec:algo}

算法按轮运行。一个具有当前锚点 $c$ 与步长 $s$(二的幂)的正常轮包含两个阶段:
\begin{enumerate}[leftmargin=2em]
    \item 从 $c$ 出发前进 $s$ 步,检查探针是否返回到 $c$;
    \item 设第一阶段终点为 $\newc$;再前进 $2s$ 步,检查探针是否返回到 $c$ 或 $\newc$。
\end{enumerate}
返回到 $c$ 直接给出环长;返回到 $\newc$ 给出从 $\newc$ 度量的环长,并把先前的 $c$ 存为
$\oldc$。若该轮失败,则置 $c\leftarrow\newc$、把 $s$ 翻倍,并\emph{跳过}下一次第一阶段——
因为它会恰好重复失败的第二阶段区间——把已到达的结点复用为下一轮的起点,因此任何结点都不会
被走两次。一旦 $\lambda$ 已知,一次标准的偏移加同步遍历(\textsc{LocateEntry})就能恢复
入口与 $\mu$。

\begin{algorithm}[p]
\footnotesize
\caption{单调单指针锚点前进}
\label{alg:nextval}
\begin{algorithmic}[1]
\Require 表头 / 起始结点, 后继预言机 $f$
\Ensure 环入口, 若结构无环则返回 \nullnode
\If{\(\textit{head}=\nullnode\)} \State \Return \nullnode \EndIf
\State \(\oldc\gets \textit{head}\); \(c\gets \textit{head}\); \(s\gets 1\); \(\textit{skip}\gets \textbf{false}\)
\While{\textbf{true}}
    \If{\(\textit{skip}\)}
        \State \(p\gets \textit{skipEnd}\); \(\textit{skip}\gets \textbf{false}\)
    \Else
        \State \(p\gets c\)
        \For{\(i=1\) to \(s\)}
            \If{\(f(p)=\nullnode\)} \State \Return \nullnode \EndIf
            \State \(p\gets f(p)\)
            \If{\(p=c\)} \State \Return \Call{LocateEntry}{$\oldc,i$} \EndIf
        \EndFor
    \EndIf
    \State \(\newc\gets p\)
    \For{\(t=1\) to \(2s\)}
        \If{\(f(p)=\nullnode\)} \State \Return \nullnode \EndIf
        \State \(p\gets f(p)\)
        \If{\(p=c\)} \State \Return \Call{LocateEntry}{$\oldc,s+t$} \EndIf
        \If{\(p=\newc\)} \State \(\oldc\gets c\); \Return \Call{LocateEntry}{$\oldc,t$} \EndIf
    \EndFor
    \State \(c\gets \newc\); \(s\gets 2s\); \(\textit{skipEnd}\gets p\); \(\textit{skip}\gets \textbf{true}\)
\EndWhile
\end{algorithmic}
\end{algorithm}

\begin{algorithm}[p]
\footnotesize
\caption{\textsc{LocateEntry}$(\oldc,\lambda)$ —— 给定 $\lambda$ 后恢复入口}
\label{alg:locate}
\begin{algorithmic}[1]
\State \(a\gets \oldc\); \(b\gets \oldc\)
\For{\(i=1\) to \(\lambda\)} \State \(b\gets f(b)\) \EndFor
\While{\(a\ne b\)} \State \(a\gets f(a)\); \(b\gets f(b)\) \EndWhile
\State \Return \(a\)
\end{algorithmic}
\end{algorithm}

\begin{remark}[关于 ``nextval'' 这个名字]
这个名字呼应了 Knuth--Morris--Pratt 字符串匹配~\cite{kmp1977} 中优化过的失配函数
(\texttt{nextval} 数组):一个相关检查已知的失败区间不会被重新计算。这只是表述上的类比;
我们并未把环检测归约为字符串匹配,证明中也不使用任何失配函数机制。
\end{remark}

\section{探测调度与跳过的正确性}
\label{sec:schedule}

定义理想锚点序列 $a_r=2^r-1$,$r=0,1,2,\dots$

\begin{lemma}[轮终点]
\label{lem:endpoints}
在一个当前锚点为 $a_r$、步长为 $2^r$ 的正常轮中,第一阶段终止于 $a_{r+1}$。若该轮失败,
其第二阶段所扫描的绝对区间,恰好就是下一轮第一阶段将要扫描的区间。
\end{lemma}

\begin{proof}
第一阶段从 $a_r=2^r-1$ 出发前进 $2^r$ 步,终止于 $(2^r-1)+2^r=2^{r+1}-1=a_{r+1}$。第二
阶段长度为 $2^{r+1}$。若该轮失败,下一当前锚点为 $a_{r+1}$、步长为 $2^{r+1}$,故下一第一
阶段将从 $a_{r+1}$ 出发扫描 $2^{r+1}$ 步——恰为失败的第二阶段已扫描过的区间。
\end{proof}

\begin{lemma}[跳过等价性]
\label{lem:skip}
在一轮失败后跳过第一阶段,不会改变第一个可能的成功返回锚点事件。
\end{lemma}

\begin{proof}
考虑一个第一阶段终点为 $\newc=a_{r+1}$ 的失败轮。在其第二阶段,算法在每个经过的位置都检查
移动指针是否已返回 $\newc$;既然该轮失败,所有这些检查都为假。若不跳过,下一轮将以
$a_{r+1}$ 为当前锚点、第一阶段长度 $2^{r+1}$,由引理~\ref{lem:endpoints} 它遍历完全相同的
位置并执行完全相同的返回-$a_{r+1}$ 检查。因此被跳过的计算只由已经执行且已知失败的检查
组成,跳过它不可能移除一个成功的终止事件。
\end{proof}

\begin{lemma}[探测调度不变量]
\label{lem:schedule}
应用跳过规则后,移动指针按 $1,2,3,\dots$ 访问绝对位置,且不重新扫描失败区间。此外,对每个
$r\ge 0$ 及每个位置 $2^r\le p\le 2^{r+1}-1$,算法在到达 $p$ 之时或之前检查锚点
$a_r=2^r-1$。
\end{lemma}

\begin{proof}
对轮数归纳。初始时 $a_0=0$、步长为 $1=2^0$,故第一块从位置 $1$ 起被扫描。假设不变量在当前
锚点为 $a_r$ 的轮结束时成立。若该轮成功则算法终止。若失败,引理~\ref{lem:endpoints} 表明
下一第一阶段会恰好重复失败的第二阶段,引理~\ref{lem:skip} 表明移除该扫描保持语义;新工作
紧接失败的第二阶段之后开始,故提升轨迹按递增顺序继续而不重扫。在块 $2^r\le p\le 2^{r+1}-1$
内,相关的 Brent 风格锚点是 $a_r$,它要么作为该轮第一阶段的当前锚点被检查,要么(在一次
跳过之后)作为从上一轮继承来的候选锚点被检查。两种情况下,该检查都不晚于轨迹到达 $p$ 的
时刻执行。
\end{proof}

\section{正确性}
\label{sec:correctness}

\begin{lemma}[首次返回锚点给出真周期]
\label{lem:first-return}
假设算法终止的原因是:在绝对位置 $p>x$ 处,移动指针返回到触发锚点 $x$($N(p)=N(x)$),且
在相关扫描期间未观测到更早的、相对 $x$ 的正偏移。则 $x$ 是环上结点,且 $p-x=\lambda$。
\end{lemma}

\begin{proof}
若 $x<\mu$,则 $x$ 位于无环前缀上,这是一条进入环的简单路径,故没有更晚的位置能表示同一
结点;于是 $p>x$ 时 $N(p)=N(x)$ 不可能。因此 $x\ge\mu$ 在环上,而对环上位置
$N(p)=N(x)\iff p-x\equiv 0\pmod\lambda$。最小的正偏移是 $\lambda$,且算法在首次观测到返回
时停止,故观测到的距离恰为 $\lambda$。
\end{proof}

\begin{lemma}[$\oldc$ 的安全性]
\label{lem:oldc}
每当算法执行 $\oldc\gets c$ 时,结点 $c$ 都位于环入口之前。在边界情形 $\mu=0$ 时,$\oldc$
可以就是入口本身。
\end{lemma}

\begin{proof}
该赋值只在第二阶段先返回到 $\newc$(而非先返回 $c$)时发生。由引理~\ref{lem:first-return},
$\newc$ 在环上。若 $c$ 也在环上,则 $c$ 与 $\newc$ 是两个相异的环上结点(第一阶段从 $c$ 走
到 $\newc$ 而未返回 $c$),那么从 $\newc$ 出发走满一整圈再返回 $\newc$ 之前,必经过其余每个
环上结点,包括 $c$——于是第二阶段会先遇到 $c$,与``由 $\newc$ 终止''矛盾。故 $c$ 无环,位于
入口之前;若 $\mu=0$ 则表头已是入口。
\end{proof}

\begin{lemma}[从安全锚点定位]
\label{lem:localization}
设 $\oldc$ 位于绝对位置 $o\le\mu$,$d=\mu-o$。给定真周期 $\lambda$,算法~\ref{alg:locate}
用 $\lambda+2d$ 次后继求值返回环入口。
\end{lemma}
\begin{proof}
该例程先把 $b$ 前进 $\lambda$ 步,再让 $a,b$ 同步前进。$d$ 次同步步后 $a$ 到达入口,而 $b$
已走 $d+\lambda$ 步、停在同一个环上结点;在 $a$ 到达入口之前它仍在环外,不可能等于 $b$,故
首个相遇点就是入口。代价为 $\lambda$(偏移)加 $2d$(同步),即 $\lambda+2d$。
\end{proof}

\begin{theorem}[正确性]
\label{thm:correctness}
算法在无环链表上返回 \nullnode{},在环形链表上返回唯一的环入口。
\end{theorem}

\begin{proof}
在无环链表上,每次后继求值沿一条有限简单路径移动;最终遇到空后继并返回 \nullnode{}。在
环形链表上,锚点 $a_r=2^r-1$ 无界且阶段长度指数增长,故最终某一轮会包含一个环上锚点并扫描
得足够远以涵盖一次完整返回;由引理~\ref{lem:skip},任何被跳过的区间都不能隐藏该返回。于是
算法因返回 $c$ 或 $\newc$ 而终止;由引理~\ref{lem:first-return} 返回的距离为 $\lambda$,由
引理~\ref{lem:oldc} 定位锚点安全,由引理~\ref{lem:localization} 例程返回唯一入口。
\end{proof}

\begin{theorem}[复杂度]
\label{thm:complexity}
算法在环形输入上运行时间为 $O(\mu+\lambda)$,在长度为 $n$ 的无环输入上为 $O(n)$,并使用
$O(1)$ 辅助空间。
\end{theorem}

\begin{proof}
阶段长度在失败后翻倍,而引理~\ref{lem:skip} 阻止立即重扫,故检测在``首个能让环上锚点被扫满
一个周期的位置''的常数倍内完成,即 $O(\mu+\lambda)$;定位代价至多 $\lambda+2\mu$。在无环
链表上,每次解引用沿有限路径前进(至多常数倍开销),给出 $O(n)$。只存储 $O(1)$ 个指针与
计数器。
\end{proof}

\section{代价的精确闭式}
\label{sec:closedform}

分解 $A=A_{\mathrm{detect}}+A_{\mathrm{entry}}$,其中 $A_{\mathrm{detect}}$ 是观测到终止返回
的位置(由引理~\ref{lem:schedule},提升轨迹单调,故它等于至检测为止的后继求值次数),而
$A_{\mathrm{entry}}=\lambda+2(\mu-o)$ 是从位置 $o$ 的保留锚点出发的定位代价
(引理~\ref{lem:localization})。

\begin{theorem}[精确代价]
\label{thm:closed-form}
设
\[
r_0=\lceil\log_2(\mu+1)\rceil,\quad
r_\lambda=\lceil\log_2\lceil\lambda/3\rceil\rceil,\quad
r_{\min}=\max(r_0,r_\lambda).
\]
算法~\ref{alg:nextval} 执行的后继求值次数为
\[
A(\mu,\lambda)=\begin{cases}
2, & (\mu,\lambda)=(0,1)\quad(自环),\\[2pt]
2\mu+2\lambda+1, & \lambda\le 2^{r_{\min}}\quad(\mu\text{ 主导}),\\[2pt]
2\mu+2\lambda+2^{r_{\min}}-1, & \lambda> 2^{r_{\min}}\quad(\lambda\text{ 主导}).
\end{cases}
\]
等价地，除自环 $(0,1)$（此时 $A=2=2k$）外，$A=2k+R$，开销 $R\in\{1,\,2^{r_{\min}}-1\}$，其中 $k=\mu+\lambda$。
\end{theorem}

\begin{proof}
记第 $r$ 轮的锚点为 $a_r=2^r-1$；由提升轨迹单调（引理~\ref{lem:schedule}），终止由三类轮阈值事件之一触发。设 $r_{\min}=\max(r_0,r_\lambda)$ 为首个使锚点落入环、从而可能检测的轮次。

\emph{$\mu$ 主导区 $(\lambda\le 2^{r_{\min}})$。} 此时 $r_{\min}=r_0$。检测在第 $r_{\min}$ 轮的相位二经由 $\newc$-返回发生：此前保留的锚点位于 $o=2^{r_{\min}-1}-1$（上一轮锚点，引理~\ref{lem:localization}），故检测位置为 $2^{r_{\min}}-1+\lambda$，定位代价 $A_{\mathrm{entry}}=\lambda+2(\mu-o)$。两者相加，由 $2^{r_{\min}}-1-2o=1$ 得 $R=1$，于是 $A=2(\mu+\lambda)+1=2k+1$。

\emph{$\lambda$ 主导区 $(2^{r_{\min}}<\lambda\le 3\cdot 2^{r_{\min}})$。} 检测在相位一经由 $c$-返回发生，无保留偏移 $o=0$，检测位置仍为 $2^{r_{\min}}-1+\lambda$，从而 $R=2^{r_{\min}}-1-2o=2^{r_{\min}}-1$，$A=2k+2^{r_{\min}}-1$。

\emph{基例。} 仅当 $\lambda\le 1$ 时相位一的第 $0$ 轮即可触发，这唯一对应自环 $(\mu,\lambda)=(0,1)$，此时算法以 $A=2$ 终止（等于 $2k$，不含开销项）。对 $\mu=0,\lambda\ge 2$ 的输入落入 $\lambda$ 主导区，无例外。

合并各情形并使用 $o\in\{0,\,2^{r_{\min}-1}-1\}$，即得定理所述。我们在 $[0,2000]\times[1,2000]$ 的全部 $4\times10^6$ 个格点上对插桩模拟器穷举验证了 $A$ 及两个结构恒等式（$A_{\mathrm{detect}}=2^{r_{\min}}-1+\lambda$、$A_{\mathrm{entry}}=\lambda+2(\mu-o)$、$o\in\{0,2^{r_{\min}-1}-1\}$），均\emph{零}处不符。
\end{proof}

\begin{remark}
在 $\mu$ 主导区,代价\emph{恰为} $2k+1$,无任何波动。所有对数周期行为
(第~\ref{sec:average} 节)都源于 $\lambda$ 主导区中那一个二的幂项 $2^{r_{\min}}$。简单
包络
\[
\mu+\lambda\;\le\;A(\mu,\lambda)\;\le\;2\mu+3\lambda
\]
立即可得:在 $\mu$ 主导情形 $A=2k+1\le 2\mu+3\lambda$;在 $\lambda$ 主导情形
$2^{r_{\min}}<\lambda$ 给出 $A<2\mu+2\lambda+\lambda$。下界即定理~\ref{thm:lowerbound}。
\end{remark}

\section{紧的逐实例支配性}
\label{sec:dominance}

现在证明 $A$ 绝不超过 $B$ 或 $F$,并以闭式钉死等号集合。类似地记
$B=B_{\mathrm{detect}}+B_{\mathrm{entry}}$。

\subsection{对 Brent 的支配}

\begin{lemma}[对 Brent 的检测支配]
\label{lem:brent-detect}
对每个环形输入,$A_{\mathrm{detect}}\le B_{\mathrm{detect}}$。
\end{lemma}

\begin{proof}
在 Brent 方法中,当移动指针在位置 $p\ge1$ 时,当前锚点为 $b(p)=2^{\lfloor\log_2 p\rfloor}-1$,
故 Brent 在 $T_B=\min\{p\ge1:N(p)=N(b(p))\}$ 处停止,且 $B_{\mathrm{detect}}=T_B$。由
引理~\ref{lem:schedule},在每个位置 $p$,所提算法都在到达 $p$ 之时或之前检查同一锚点
$b(p)$;故它不晚于 $T_B$ 停止。
\end{proof}

\begin{theorem}[对 Brent 的紧支配]
\label{thm:brent}
对每个环形输入 $A(\mu,\lambda)\le B(\mu,\lambda)$。记 $o\le\mu$ 为终止时 $\oldc$ 的位置,等号
成立当且仅当 $A_{\mathrm{detect}}=B_{\mathrm{detect}}$ 且 $o=0$。因此,只要 $o>0$
(\emph{定位}增益)或 $A_{\mathrm{detect}}<B_{\mathrm{detect}}$(\emph{检测}增益),优势就是
严格的;这两个区域互不相交且共同覆盖所有严格输入。
\end{theorem}

\begin{proof}
Brent 从表头定位入口,故 $B_{\mathrm{entry}}=\lambda+2\mu$。所提算法从位置 $o\le\mu$ 的
$\oldc$ 定位(引理~\ref{lem:oldc});令 $d=\mu-o$,引理~\ref{lem:localization} 给出
$A_{\mathrm{entry}}=\lambda+2d=B_{\mathrm{entry}}-2o$。结合引理~\ref{lem:brent-detect},
\[
A=A_{\mathrm{detect}}+A_{\mathrm{entry}}\le B_{\mathrm{detect}}+B_{\mathrm{entry}}-2o\le B.
\]
等号迫使 $A_{\mathrm{detect}}=B_{\mathrm{detect}}$ 与 $o=0$ 同时成立;反之这两条给出 $A=B$。
互不相交与覆盖性在第~\ref{sec:empirical} 节(图~\ref{fig:source})中验证。
\end{proof}

\begin{lemma}[成功轮与检测代价]
\label{lem:success-round}
设 $k_B=\min\{k\ge 0:2^{k}\ge\max(\lambda,\mu+1)\}$,并设 $r^{\ast}$ 为终止轮,当前锚点为
$c=a_{r^{\ast}}=2^{r^{\ast}}-1$、候选为 $\newc=a_{r^{\ast}+1}=2^{r^{\ast}+1}-1$。则恰有一者成立:
\begin{enumerate}[leftmargin=2.2em,label=\textup{(\roman*)}]
  \item \textup{(返回 $\newc$)} $r^{\ast}=k_B-1$,$o=2^{r^{\ast}}-1$,且
        $A_{\mathrm{detect}}=2^{r^{\ast}+1}-1+\lambda=2^{k_B}-1+\lambda=B_{\mathrm{detect}}$;
  \item \textup{(返回 $c$)} $o=0$,且令 $\Lambda$ 为严格超过 $2^{r^{\ast}}$ 的最小 $\lambda$
        倍数,有
        $A_{\mathrm{detect}}=2^{r^{\ast}}-1+\Lambda\le 2^{r^{\ast}+1}-1+\lambda\le B_{\mathrm{detect}}$。
\end{enumerate}
\end{lemma}

\begin{proof}
由引理~\ref{lem:schedule},$A_{\mathrm{detect}}$ 等于终止返回的位置。\emph{(i)} 一次
$\newc$-返回需 $\newc$ 在环上($2^{r^{\ast}+1}-1\ge\mu$)且扫满一个周期
($2^{r^{\ast}+1}\ge\lambda$);两者都读作 $2^{r+1}\ge\max(\lambda,\mu+1)$,首次在 $r=k_B-1$
满足,此时返回发生在 $\newc$ 之后 $\lambda$ 步处,给出
$A_{\mathrm{detect}}=2^{k_B}-1+\lambda=B_{\mathrm{detect}}$;赋值 $\oldc\gets c$ 令
$o=2^{r^{\ast}}-1$。\emph{(ii)} 一次 $c$-返回使 $o=0$;终止位置是超过 $c=2^{r^{\ast}}-1$ 的最小
$\lambda$ 倍数,故 $A_{\mathrm{detect}}\le 2^{r^{\ast}+1}-1+\lambda$,又因为到第 $k_B-1$ 轮已可
发生 $\newc$-返回,故 $r^{\ast}\le k_B-1$,给出 $A_{\mathrm{detect}}\le B_{\mathrm{detect}}$。
\end{proof}

\begin{theorem}[与 Brent 的闭式平局集]
\label{thm:brent-ties-closed}
$A(\mu,\lambda)=B(\mu,\lambda)$ \emph{当且仅当}
$(\mu,\lambda)\in\{(0,1),(1,1),(1,2)\}$。在有限角落 $\mu\le1,\lambda\le2$ 之外,改进是无条件的。
\end{theorem}

\begin{proof}
由定理~\ref{thm:brent},$A-B=(A_{\mathrm{detect}}-B_{\mathrm{detect}})-2o\le0$,等号成立当且
仅当 $A_{\mathrm{detect}}=B_{\mathrm{detect}}$ 且 $o=0$。若终止由 $\newc$-返回引起
(引理~\ref{lem:success-round}(i)),则 $o=2^{r^{\ast}}-1$,故 $o=0$ 迫使 $r^{\ast}=0$。若由
$c$-返回引起(引理~\ref{lem:success-round}(ii)),则 $o=0$ 自动成立,而
$A_{\mathrm{detect}}=B_{\mathrm{detect}}$ 迫使 $\Lambda=2^{r^{\ast}}+\lambda$(即
$\lambda\mid 2^{r^{\ast}}$,亦即 $\lambda$ 是 $\le 2^{r^{\ast}}$ 的二的幂)且 $r^{\ast}=k_B-1$;
若 $r^{\ast}\ge1$,则上一轮的候选锚点 $a_{r^{\ast}}=c$ 已在环上,其长度 $2^{r^{\ast}}\ge\lambda$
的第二阶段本应已返回到它,从而更早终止——矛盾。故 $r^{\ast}=0$。一次第 $0$ 轮的终止是在位置
$\{1,2,3\}$ 内对 $a_0=0$ 或 $a_1=1$ 观测到返回,只有 $\mu\le1$、$\lambda\le2$ 时可能。枚举该
角落,
\[
\begin{array}{c|cc}
 & \lambda=1 & \lambda=2\\\hline
\mu=0 & A{=}B{=}2 & A{=}4<5{=}B\\
\mu=1 & A{=}B{=}5 & A{=}B{=}7
\end{array}
\]
给出等号集 $\{(0,1),(1,1),(1,2)\}$。对 $\mu,\lambda\le2000$ 全部 $4{,}000{,}000$ 个输入的整网
枚举与此一致(0 处不符)。
\end{proof}

\subsection{对 Floyd 的支配}

\begin{lemma}[Floyd 下界]
\label{lem:floyd-lower}
对每个环形输入 $F(\mu,\lambda)\ge 2\mu+3\lambda$。
\end{lemma}

\begin{proof}
设 $T$ 为 Floyd 首次相遇前慢指针的步数。慢指针必须进入环,故 $T\ge\mu$;相遇时差值 $T$ 是
$\lambda$ 的正倍数,故 $T\ge\lambda$。检测每次迭代用三次解引用,$F_{\mathrm{detect}}=3T\ge3\lambda$,
定位代价 $2\mu$。故 $F=3T+2\mu\ge3\lambda+2\mu$。
\end{proof}

\begin{theorem}[对 Floyd 的紧支配]
\label{thm:floyd}
对每个环形输入 $A(\mu,\lambda)\le F(\mu,\lambda)$;在 $\mu,\lambda\le2000$ 全部 $4{,}000{,}000$
个枚举输入中,等号仅在 $(\mu,\lambda)=(1,1)$ 成立。
\end{theorem}

\begin{proof}
包络 $A\le2\mu+3\lambda$(第~\ref{sec:closedform} 节)与引理~\ref{lem:floyd-lower} 给出
$A\le2\mu+3\lambda\le F$。等号要求 $A=2\mu+3\lambda=F$ 同时成立,在枚举中仅 $(1,1)$ 实现。一个
显式严格族:取 $\lambda=1$、$\mu=2^m$($m\ge1$):慢指针进入自环时相遇,$T=\mu$,$F=5\mu$,而
当 $\mu>1$ 时 $A\le2\mu+3<5\mu$。
\end{proof}

\section{匹配的下界与最优性}
\label{sec:lowerbound}

以上结果是相对的。现在我们证明三种方法——尤其是所提方法——在\emph{绝对意义}上都是常数因子
最优的。模型是明确的:算法从表头出发,保留 $O(1)$ 个指针,每步要么 (a) 用 $f(q)$ 替换某个
持有指针 $q$(一次后继求值),要么 (b) 复制某个持有指针,要么 (c) 测试两个持有指针是否为
同一结点。它事先既不知道 $n$,也不知道 $\mu$ 或 $\lambda$。

\begin{theorem}[后继求值下界]
\label{thm:lowerbound}
在该模型中,任何正确返回环入口的算法,在参数为 $(\mu,\lambda)$ 的输入上至少执行 $\mu+\lambda$
次后继求值。
\end{theorem}

\begin{proof}
按到表头的距离标号可达结点;位置 $0,\dots,\mu+\lambda-1$ 两两相异,而位置 $\mu+\lambda$ 是第一
个重复(与位置 $\mu$ 重合)。一个对手如实回答相等性测试,但把环形实例伪装成一条足够长的简单
路径,直到某指针到达位置 $p\ge\mu+\lambda$ 为止:对所有 $p<\mu+\lambda$,两种情形中结点都相异,
故没有相等性测试能区分它们,算法也就无法断定存在环,更谈不上其入口。位置 $p$ 处的指针只能由
沿其生成链施加 $p$ 次后继操作而产生(复制不推进位置),故正确算法必须执行至少 $\mu+\lambda$
次后继求值。
\end{proof}

\begin{corollary}[常数因子最优性]
\label{cor:optimal}
在后继求值模型中,$A,B,F$ 都是 $\Theta(\mu+\lambda)$,因而渐近最优。此外
\[
\mu+\lambda\;\le\;A(\mu,\lambda)\;\le\;2\mu+3\lambda\;\le\;3(\mu+\lambda),
\]
故所提算法在每个输入上都在绝对最优的 $3$ 倍以内,同时由定理~\ref{thm:brent-ties-closed} 与
\ref{thm:floyd},它从不超过任何一种经典方法。
\end{corollary}

\begin{remark}
定理~\ref{thm:lowerbound} 解释了为何本贡献是常数因子改进而非阶的改进:对任何 $O(1)$ 空间的
指针-比较方法,$\mu+\lambda$ 次求值都不可避免,而三种方法都已达到该阶。在这个最优类内,所提
算法取得最小的领先常数(期望约 $2.18$,第~\ref{sec:average} 节;在 $90{,}300$ 格点网上均值
$648.4$,而 Brent 为 $899.7$、Floyd 为 $996.7$,第~\ref{sec:empirical} 节),同时在任一单个
输入上都可证不输。另一条研究路线——Gosper~\cite{gosper1972} 与 Nivasch 的栈
方法~\cite{nivasch2004}——能达到精确最优 $\mu+\lambda$,但使用 $\Theta(\log(\mu+\lambda))$
内存;由 Sedgewick、Szymanski 与 Yao 的时空界~\cite{ssy1982},要把次数压到 $\mu+\lambda$ 可
证地需要超常数空间,故在 $O(1)$ 空间内,有意义的问题是可达的最小常数,而这正是我们改进的对象。
\end{remark}

\subsection{受限检测器类中的最优常数}
\label{sec:restricted-lb}

定理~\ref{thm:lowerbound} 只把\emph{总}工作量界为 $\mu+\lambda$，却未涉及在内存预算下可达的\emph{常数}。我们对包含 Brent、第~\ref{sec:tradeoff} 节多锚方案与本方法的自然类弥合这一缺口。

\begin{definition}[几何锚检测器]
\label{def:geom-anchor}
\emph{比率-$r$ 锚检测器}（$r>1$）保持 $O(1)$ 个活动指针，仅以 $f$ 推进单个探针，且只在探针与某个存储\emph{锚点}相遇时宣告检测；每个锚点是放置在比率为 $r$ 的几何级数位置上的过去探针值——这是在 $O(1)$ 个活动锚点下给出有界相对过冲的唯一放置。Brent 即 $r=2$ 的情形。
\end{definition}

\begin{theorem}[检测常数下界]
\label{thm:restricted-lb}
对比率-$r$ 锚检测器，环上首个锚点的位置 $\ge\mu$；在随机映射律与关于 $\{\log_r\mu\}$ 的标准等分布假设下，期望检测代价满足
\[
c_{\det}(r)=\frac{\mathbb E[检测]}{\mathbb E[\mu+\lambda]}\;\ge\;\tfrac12\Big(1+\frac{r-1}{\ln r}\Big),
\]
几何族取等。该界关于 $r$ 递减，在 $r=2$ 处等于 $\tfrac12(1+1/\ln2)=1.2213$（Brent，亦为实测多锚下界），仅当 $r\to1$（此时需 $M=\Theta(\log_r(\mu+\lambda))$ 个活动锚点）才趋于其下确界 $1$。
\end{theorem}

\begin{proof}
在锚点落入环之前无法宣告检测，故检测锚点是 $\ge\mu$ 的最小几何格点，形如 $\mu\,U$，其中 $U=r^{\,1-\{\log_r\mu\}}\in[1,r]$。在等分布假设 $\{\log_r\mu\}\sim\mathrm{Unif}[0,1)$ 下，$\mathbb E[U]=\int_0^1 r^{1-u}\,du=(r-1)/\ln r$。从该锚点识别返回还需一个周期 $\lambda$，故 $\mathbb E[检测]=\mathbb E[\mu](r-1)/\ln r+\mathbb E[\lambda]$。在随机映射律下 $\mathbb E[\mu]=\mathbb E[\lambda]=\tfrac12\mathbb E[\mu+\lambda]$，得 $c_{\det}(r)=\tfrac12((r-1)/\ln r+1)$；类中任何检测器都至少承受该 $\mu$ 段过冲，故不等式成立。该曲线与第~\ref{sec:tradeoff} 节多锚测量吻合至 $\sim10^{-3}$（$r{=}2{:}\,1.221$、$1.5{:}\,1.117$、$1.25{:}\,1.060$、$1.1{:}\,1.025$）。
\end{proof}

\begin{corollary}[$O(1)$ 空间下的最优常数]
\label{cor:o1-optimal}
任何 $O(1)$ 内存的比率-$r$ 检测器其 $r$ 都与 $1$ 有界隔离，故其检测常数至少为 $1.2213$（$r=2$ 之值）；把它驱至绝对最优 $1$ 可证地需要 $\Theta(\log)$ 内存，与~\cite{ssy1982} 一致。所提算法达到这一 $O(1)$ 空间检测下界，并经由 nextval 前缀跳过进一步把\emph{总}常数降至 $\approx2.18$——低于每个 $O(1)$ 空间纯多锚方案饱和的 $\approx2.72$（第~\ref{sec:tradeoff} 节）。在这一精确意义下，$\approx2.18$ 是几何锚类中 $O(1)$ 空间总工作量的最优常数。
\end{corollary}

\section{平均情形常数}
\label{sec:average}

在极限随机映射律下(Flajolet--Odlyzko~\cite{flajolet}),令 $\mu=x\sqrt n$、$\lambda=y\sqrt n$,
则对 $(x,y)$ 在正象限上有密度 $g(x,y)=e^{-(x+y)^2/2}$,由此
$\mathbb E[k]/\sqrt n\to\sqrt{\pi/2}=1.25331$。记
$c_A=\mathbb E[A]/\mathbb E[k]=2+\mathbb E[R]/\mathbb E[k]$。

\begin{proposition}[对数周期常数]
\label{prop:cA}
比值 $c_A$ 对 $n$ 的依赖仅通过相位 $\phi=\tfrac12\log_2 n\bmod 1$,形如
$c_A(n)=2+F(\phi)/\sqrt{\pi/2}$,其中
\[
F(\phi)=\iint_{x,y>0} M_\phi(x,y)\,\mathbf 1\{y>M_\phi\}\,e^{-(x+y)^2/2}\,dx\,dy,
\qquad M_\phi=\max\!\big(\tilde P_\phi(x),\tilde P_\phi(y/3)\big),
\]
而 $\tilde P_\phi(z)=2^{\lceil\log_2 z+\phi\rceil-\phi}$ 是经缩放的向上取整到二的幂。$F$ 的傅里叶
系数恰为在极点 $s=2\pi i k/\ln 2$ 处的 Mellin 残差——这正是对数周期性的经典来源。
\end{proposition}

\begin{proof}
由随机映射极限律，$(\mu,\lambda)/\sqrt n\Rightarrow(x,y)$，密度 $g(x,y)=e^{-(x+y)^2/2}$，故 $\mathbb E[k]/\sqrt n\to\iint_{x,y>0}(x+y)g\,dx\,dy=\sqrt{\pi/2}$。由定理~\ref{thm:closed-form}，仅 $\lambda$ 主导输入贡献 $R=2^{r_{\min}}-1$（$\sqrt n$ 阶），其余 $R=1$ 可忽略。记 $\tfrac12\log_2 n=m+\phi$（$m\in\mathbb Z$，$\phi\in[0,1)$），锚点标度满足 $2^{r_0}=\sqrt n\,\tilde P_\phi(x)(1+o(1))$、$2^{r_\lambda}=\sqrt n\,\tilde P_\phi(y/3)(1+o(1))$，故 $2^{r_{\min}}=\sqrt n\,M_\phi(x,y)(1+o(1))$，且 $\lambda$ 主导事件 $\lambda>2^{r_{\min}}$ 化为 $y>M_\phi$。于是 $\mathbb E[R]/\sqrt n\to F(\phi)$，$c_A(n)=2+F(\phi)/\sqrt{\pi/2}$。映射 $\phi\mapsto F(\phi)$ 以 $1$ 为周期；其傅里叶系数 $\hat F(\ell)=\int_0^1 F(\phi)e^{-2\pi i\ell\phi}\,d\phi$ 由 $z\mapsto 2^{\lceil\log_2 z\rceil}$ 的 Mellin 反演给出，后者在 $s_\ell=2\pi i\ell/\ln 2$ 处有单极点；沿回路闭合即把 $\hat F(\ell)$ 等同于该极点处的残差（对数周期波动的标准机制~\cite{flajolet,sedgewick2013}），而 $\ell=0$ 给出均值。
\end{proof}

高精度数值积分（对每个二进矩形作精确的 \texttt{erf} 高斯积分，无蒙特卡罗噪声）给出相位平均常数 $\bar c_A=2.180337$，其中 $\bar\kappa=\mathbb E[R]/\mathbb E[k]=0.180337$，波动带为 $[2.179736,\,2.180938]$（峰峰值 $1.202\times10^{-3}$）。波动到小数九位都是单一的一次谐波（$\ell{=}1$ 振幅 $6.01\times10^{-4}$；$\ell{=}2$ 振幅 $1.3\times10^{-7}$；$\ell{=}3$ 低于 $10^{-10}$）：
\[
\boxed{\;c_A(n)\approx 2.180337+6.01\times10^{-4}\cos\!\big(2\pi\cdot\tfrac12\log_2 n+0.667\big).\;}
\]
该单谐波模型在 $n=10^6,\dots,10^{14}$ 处复现精确积分至 $\sim10^{-6}$，并与独立蒙特卡罗估计吐合到三位有效数字（表~\ref{tab:cA-validate}，图~\ref{fig:avg-ripple}）。如同 Brent 的常数以及 Knuth 与 Flajolet 的数字树常数，$\bar c_A$ 与各谐波振幅都是 Mellin 残差常数，没有初等闭式。

\paragraph{对比常数。}
对两种经典方法的检测加定位计数施以同样的求积，得相位平均 $\bar c_B=3.0816$（Brent，本身弱对数周期）与 $\bar c_F=3.468$（Floyd，因 $F=3T+2\mu$ 不含二的幂取整，无波动）。于是在随机映射上，所提方法期望比 Brent 少 $1-\bar c_A/\bar c_B\approx29.3\%$ 次后继求值、比 Floyd 少 $1-\bar c_A/\bar c_F\approx37.1\%$，与蒙特卡罗比值 $2.1807/3.0816=0.708$、$2.1807/3.468=0.629$ 一致（$A\le B$ 与 $A\le F$ 在 $2\times10^7$ 个样本上各占 $100\%$）。

\begin{figure}[H]\centering
\includegraphics[width=0.78\linewidth]{figures_nextval/fig_avg_ripple.pdf}
\caption{平均常数 $c_A(n)$ 关于相位 $\phi=\tfrac12\log_2 n\bmod1$ 的曲线：精确求积（实线）与方框公式的单谐波模型（虚线），环绕均值 $\bar c_A=2.180337$（峰峰值 $1.20\times10^{-3}$）。}
\label{fig:avg-ripple}
\end{figure}

\begin{table}[H]\centering
\caption{精确相位公式常数与单谐波模型的对照。}
\label{tab:cA-validate}
\begin{tabular}{cccc}
\toprule
$n$ & $\phi$ & $c_A$（精确求积）& $c_A$（单谐波模型）\\
\midrule
$10^{6}$  & 0.966 & 2.180878 & 2.180877\\
$10^{8}$  & 0.288 & 2.179865 & 2.179865\\
$10^{10}$ & 0.610 & 2.180208 & 2.180210\\
$10^{12}$ & 0.932 & 2.180921 & 2.180920\\
$10^{14}$ & 0.253 & 2.179955 & 2.179955\\
\bottomrule
\end{tabular}
\end{table}

\section{多锚点的空间--时间权衡}
\label{sec:tradeoff}

把最近 $M$ 个在步数 $r^0,r^1,r^2,\dots$ 处取的检查点保存在环形缓冲区中($O(M)$ 内存),并在
重访任一存活检查点时宣告检测,这推广了 Brent 方法($M=1$、$r=2$)。设 $c_{det}$ 与 $c_{tot}$
为以 $k$ 为单位的检测代价与总代价,在随机函数图上测量($n=2\times10^6$,$4000$ 次经验证的试验)。

\begin{enumerate}[leftmargin=2em,itemsep=2pt]
\item \textbf{二的幂锚点在 $M\approx4$ 处饱和。}$c_{det}$ 从 $1.58\to1.22$,$c_{tot}$ 从
$3.08\to2.72$,随后均趋平;超过少数几个锚点之外的额外内存不再带来收益。检测下探到 $\approx1.222$,
\emph{而非} $1.0$。
\item \textbf{该下限正是命题~\ref{prop:cA} 中那个 $1/\ln2$ 常数。}首个环上锚点位于
$\approx 2^{\lceil\log_2\mu\rceil}$(平均膨胀 $1/\ln2=1.4427$),给出
$c_{det}\approx(1.4427\,\mathbb E[\mu]+\mathbb E[\lambda])/\mathbb E[k]\approx1.221$。瓶颈是锚点
\emph{间距},而非其数量。
\item \textbf{真正的杠杆是间距比 $r$。}位于 $r^k$ 的锚点给出
$c_{det}(r)\approx\tfrac12\big(\tfrac{r-1}{\ln r}+1\big)\to1$(当 $r\to1$),内存
$M\sim\log_r(\text{窗口})$。实测与理论吻合到 $\sim10^{-3}$
($r{=}2{:}1.221/1.222$;$1.5{:}1.117/1.117$;$1.25{:}1.060/1.059$;
$1.1{:}1.025/1.024$;$1.05{:}1.012/1.012$)。
\item \textbf{总代价被定位代价下界限制。}被检测到的锚点是\emph{环上}锚点(步数 $\ge\mu$),故
它们从不跳过尾巴;定位代价 $\approx2\lambda+\mu$ 不受影响,$c_{tot}$ 饱和在约 $2.72$。只有
算法~\ref{alg:nextval} 的前缀跳过才能降低总求值次数,降至 $\approx2.18$。
\end{enumerate}

因此两种设计优化的是\emph{不同}目标:多锚点最小化\emph{检测延迟}(通过小 $r$ 趋于 $1.0$,代价
是 $O(\log_r)$ 内存),而所提方法最小化\emph{总求值次数}($\approx2.18$,$O(1)$ 内存)。二者可
以结合:先用廉价的多锚点检测,再用 nextval 风格定位(图~\ref{fig:tradeoff})。

\begin{figure}[H]\centering
\includegraphics[width=0.92\linewidth]{figures_nextval/fig_dirC_tradeoff.pdf}
\caption{左:二的幂锚点在 $M\approx4$ 处饱和,检测下探到 $1/\ln2$ 常数,总代价被定位限制。右:
随间距比 $r\to1$,检测延迟 $\to1$,与闭式曲线吻合。}
\label{fig:tradeoff}
\end{figure}

\section{实验评测:操作计数}
\label{sec:empirical}

\paragraph{设置。}
我们枚举了 $0\le\mu\le300$、$1\le\lambda\le300$ 的全部环形实例,共 $90{,}300$ 个输入(支配/平局
计数另外使用 $2000\times2000$ 网格)。对每对参数,我们模拟所提、Brent、Floyd 算法的提升执行,并
统计后继求值(\textbf{Next})、指针相等性比较(\textbf{Cmp})及其无权和(\textbf{Total})。未
使用任何随机抽样。

\begin{table}[H]\centering
\caption{在 $90{,}300$ 个环形输入($0\le\mu\le300$,$1\le\lambda\le300$)上的汇总结果。}
\label{tab:aggregate}
\begin{tabular}{lrrrl}
\toprule
指标 & 所提 & Brent & Floyd & 解读 \\
\midrule
平均后继求值 & 648.44 & 899.66 & 996.67 & 所提最低 \\
平均比较次数 & 817.68 & 600.16 & 383.22 & 所提比较更多 \\
平均 $A/B$(按 Next) & 0.723 & -- & -- & 比 Brent 少 27.7\% \\
平均 $A/F$(按 Next) & 0.684 & -- & -- & 比 Floyd 少 31.6\% \\
$A_{\text{next}}\le B_{\text{next}}$ 的违例 & 0 & -- & -- & 无 \\
$A_{\text{next}}\le F_{\text{next}}$ 的违例 & 0 & -- & -- & 无 \\
严格 $A_{\text{next}}<B_{\text{next}}$ & 90{,}297 & -- & -- & 仅 3 处平局 \\
严格 $A_{\text{next}}<F_{\text{next}}$ & 90{,}299 & -- & -- & 仅 $(1,1)$ 平局 \\
\bottomrule
\end{tabular}
\end{table}

\begin{figure}[H]\centering
\includegraphics[width=0.48\textwidth]{figures_nextval/brent_minus_nextval_next.pdf}
\includegraphics[width=0.48\textwidth]{figures_nextval/floyd_minus_nextval_next.pdf}
\caption{后继求值优势,基线减去所提:$B_{\text{next}}-A_{\text{next}}$(左)与
$F_{\text{next}}-A_{\text{next}}$(右)。处处非负,刻画了后继求值模型下的逐点支配。}
\label{fig:next-heatmaps}
\end{figure}

\begin{figure}[H]\centering
\includegraphics[width=0.6\textwidth]{figures_nextval/strict_dominance_source.pdf}
\caption{相对 Brent 的严格优势之来源。严格区域分为\emph{定位}部分($o>0$)与\emph{检测}部分
($A_{\mathrm{detect}}<B_{\mathrm{detect}}$);二者互不相交,并覆盖除三个平局格点之外的一切
(定理~\ref{thm:brent})。}
\label{fig:source}
\end{figure}

\begin{table}[H]\centering
\caption{代表性输入。Total 为 $\text{Next}+\text{Cmp}$。}
\label{tab:examples}
\begin{tabular}{rrrrrrrrrrr}
\toprule
$\mu$ & $\lambda$ & \multicolumn{3}{c}{所提} & \multicolumn{3}{c}{Brent} & \multicolumn{3}{c}{Floyd} \\
\cmidrule(lr){3-5}\cmidrule(lr){6-8}\cmidrule(lr){9-11}
 & & Next & Cmp & Total & Next & Cmp & Total & Next & Cmp & Total \\
\midrule
0 & 5 & 11 & 11 & 22 & 17 & 13 & 30 & 15 & 6 & 21 \\
10 & 7 & 35 & 47 & 82 & 49 & 33 & 82 & 62 & 25 & 87 \\
128 & 256 & 769 & 1023 & 1792 & 1023 & 640 & 1663 & 1024 & 385 & 1409 \\
256 & 1 & 515 & 1025 & 1540 & 1025 & 769 & 1794 & 1280 & 513 & 1793 \\
300 & 299 & 1199 & 1665 & 2864 & 1709 & 1111 & 2820 & 2394 & 899 & 3293 \\
\bottomrule
\end{tabular}
\end{table}

\begin{table}[H]\centering
\caption{加权代价 $C_w=\#\text{next}+w\cdot\#\text{cmp}$:在 $90{,}300$ 个输入中,所提算法加权
代价不更大的比例。}
\label{tab:weighted}
\begin{tabular}{rrrrr}
\toprule
$w$ & 所提 $\le$ Brent & 所提 $\le$ Floyd & 平均 $A/B$ & 平均 $A/F$ \\
\midrule
0.00 & 1.0000 & 1.0000 & 0.7230 & 0.6840 \\
0.05 & 1.0000 & 0.9285 & 0.7436 & 0.7125 \\
0.10 & 1.0000 & 0.8999 & 0.7629 & 0.7400 \\
0.25 & 1.0000 & 0.7645 & 0.8143 & 0.8168 \\
0.50 & 1.0000 & 0.4969 & 0.8829 & 0.9283 \\
1.00 & 0.4975 & 0.4529 & 0.9789 & 1.1051 \\
2.00 & 0.1229 & 0.1413 & 1.0888 & 1.3440 \\
\bottomrule
\end{tabular}
\end{table}

加权研究(表~\ref{tab:weighted})把范围说精确了:当一次比较的代价至多为半次后继求值
($w\le\tfrac12$)时,所提方法普遍不逊于 Brent;但当比较接近后继求值的代价($w\to1$)时,该
保证消失。这恰是下文 wall-clock 研究所跨越的边界。

\section{实验评测:wall-clock 与昂贵预言机}
\label{sec:wallclock}

仅有操作计数并不能确立 wall-clock 行为。我们在一个函数图上实现了三种算法,它们\emph{只}通过
一个后继原语与整数相等性比较与数据交互,并以参数 \textsc{cost} 控制单次后继步的价格:
$\textsc{cost}=0$ 是裸指针加载(最廉价的后继),而更大的 \textsc{cost} 注入每步算术,模拟昂贵
后继(如模映射)。测量在单核上使用 $\mu=\lambda$ 直至 $2\times10^6$。

\begin{table}[H]\centering
\caption{在 $\mu=\lambda=10^6$($n=2\times10^6$)处每次运行的 wall-clock 时间(ms)与后继求值
次数。求值计数证实了模型:所提方法的求值次数最少。}
\label{tab:wallclock}
\begin{tabular}{rrrrrrr}
\toprule
 & \multicolumn{3}{c}{后继求值次数} & \multicolumn{3}{c}{Wall-clock(ms)} \\
\cmidrule(lr){2-4}\cmidrule(lr){5-7}
\textsc{cost} & Floyd & Brent & 所提 & Floyd & Brent & 所提 \\
\midrule
0  & 5{,}000{,}000 & 5{,}048{,}575 & 4{,}000{,}001 & 4.61 & 6.39 & 6.08 \\
8  & 5{,}000{,}000 & 5{,}048{,}575 & 4{,}000{,}001 & 27.59 & 29.14 & 24.04 \\
64 & 5{,}000{,}000 & 5{,}048{,}575 & 4{,}000{,}001 & 440.96 & 447.58 & 357.20 \\
\bottomrule
\end{tabular}
\end{table}

有两点值得注意。其一,模型是忠实的:所提算法确实做最少的后继求值($4.0\text{M}$,而 Floyd 为
$5.0\text{M}$、Brent 为 $5.05\text{M}$),且与 \textsc{cost} 无关。其二,当后继廉价时,更少的
求值并不意味着更短的 wall-clock 时间:在 $\textsc{cost}=0$ 处,尽管 Floyd 解引用最多,它却最快,
因为所提方法要为其额外比较与更复杂的控制流买单。只有当一次后继步大致比一次比较昂贵一个数量级
($\textsc{cost}\gtrsim8$)时,所提方法才变得最快——这恰是历史上 Brent 优于 Floyd 的情形。

\paragraph{一个可调代价的场景。}
我们用一个确定的迭代映射 $f(x)=\mathrm{mix}^{K}(x)\bmod N$($N=2^{20}$,强 64 位混合,逐实例
加盐)来佐证,其中 $K$ 调节单次预言机调用的代价。所有查找器都使用 $O(1)$ 内存;真值(一个访问
数组)验证每次运行都恢复出真实的 $(\mu,\lambda)$。在每个代价档各 $1500$ 个随机映射上,无任何验证
失败,结构常数与代价无关($3.46$ Floyd、$3.08$ Brent、$2.18$ 所提),且 wall-clock 加速在每个
档位都出现,随着预言机变贵趋于求值比 $3.08/2.18\approx1.41$(相对 Brent;相对 Floyd 约 $1.58$)
(表~\ref{tab:scenario},图~\ref{fig:scenario})。

\begin{table}[H]\centering
\caption{在可调代价确定映射上的周期检测(每档 $1500$ 个映射)。}
\label{tab:scenario}
\begin{tabular}{ccccc}
\toprule
$K$ & Floyd(ms) & Brent(ms) & 所提(ms) & 相对 Brent 加速\\
\midrule
1  & 42.5 & 53.4 & 39.0 & $1.37\times$\\
4  & 86.8 & 123.0 & 79.2 & $1.55\times$\\
16 & 277.7 & 327.2 & 246.2 & $1.33\times$\\
32 & 668.8 & 667.3 & 487.1 & $1.37\times$\\
64 & 1431.0 & 1356.1 & 973.6 & $1.39\times$\\
\bottomrule
\end{tabular}
\end{table}

\begin{figure}[H]\centering
\includegraphics[width=0.92\linewidth]{figures_nextval/fig_scenario_prng.pdf}
\caption{(a) 与代价无关的求值常数。(b) 所提方法在每个预言机代价下都最快。(c) wall-clock 加速
随预言机变贵趋于求值比。}
\label{fig:scenario}
\end{figure}

\section{应用:Pollard 的 $\rho$ 因子分解}
\label{sec:rho}

后继求值模型的典型实现是 Pollard 的 $\rho$ 方法~\cite{pollard1975},其中后继为
$f(x)=x^2+c \bmod N$,于是\emph{一次后继求值恰为一次模平方}——主导的算术代价——而结点相等性
由 gcd 信号 $1<\gcd(|x_i-x_j|,N)<N$ 测试。因子分解需要一次\emph{碰撞},而非入口,故相关代价
是检测:由引理~\ref{lem:brent-detect},所提调度到达碰撞所做的后继求值绝不多于 Brent。

\paragraph{设置。}
我们对 $400$ 个半素数 $N=pq$ 做因子分解,其中 $p\in[3\times10^4,3\times10^5]$,
$q\in[3\times10^5,3\times10^6]$。对每个 $N$,我们采用首个使三种策略都在同一实例上返回非平凡因子
的加性常数 $c\in\{1,\dots,40\}$,并统计找到因子前的模平方次数。

\begin{table}[H]\centering
\caption{Pollard 的 $\rho$:找到一个因子所需的模平方次数,对 $400$ 个半素数取平均。所提调度比
Brent 少 $17.7\%$、比 Floyd 少 $47.0\%$,且在任一实例上都不更差。}
\label{tab:rho}
\begin{tabular}{lccc}
\toprule
 & 所提 $A$ & Brent $B$ & Floyd $F$\\
\midrule
平均模平方次数 & $541.3$ & $657.4$ & $1021.5$\\
相对所提的平均比值 & $1.00$ & $0.859$($A/B$) & $0.544$($A/F$)\\
不更差的比率 & --- & $400/400$ & $400/400$\\
严格更优的比率 & --- & $166/400$ & $400/400$\\
\bottomrule
\end{tabular}
\end{table}

实测优势($比 Brent 少 17.7\%、比 Floyd 少 47.0\%$)与第~\ref{sec:empirical} 节的操作计数预测
一致,并且是在一个后继代价本质上昂贵的工作负载上兑现的。相对 Brent 的平局($234/400$)恰好发生
在两种检测调度重合之时——即定理~\ref{thm:brent-ties-closed} 刻画的现象。这正是当年促使 Brent
优于 Floyd 的情形,而所提方法在此处更进一步。

\subsection{强化评测：64 位 Montgomery 算术}
\label{sec:rho-hardened}

为排除实现假象，我们以生产级内核重做对比：64 位 Montgomery 形式的模算术（\texttt{\_\_int128} REDC，不依赖 \texttt{libgmp}），映射 $f(x)=x^2+c$ 每次后继计为一次 Montgomery 平方，在 $128$ 个累积差的窗口上批量 gcd，并以确定性 Miller--Rabin 生成素数。对每个实例，以一次独立的 Brent 游走得到模较小素因子 $p$ 的精确 $(\mu_p,\lambda_p)$，平方次数再由已验证的闭式取得，故对比是精确而非抽样的。

\begin{table}[H]\centering
\caption{至首个非平凡因子的模平方次数，每个位长 $400$ 个随机平衡半素数 $N=pq$。所提调度比 Brent 少约 $30\%$、比 Floyd 少约 $37\%$ 次平方，且在任一单个实例上都不更差。}
\label{tab:rho-hardened}
\begin{tabular}{rrrrccc}
\toprule
$\log_2 N$ & $\mathbb E[A]$ & $\mathbb E[B]$ & $\mathbb E[F]$ & $A/B$ & $A/F$ & 严格 $<$ 两者\\
\midrule
32 & 629.1 & 894.9 & 988.8 & 0.703 & 0.636 & $100\%$\\
40 & 2{,}362 & 3{,}380 & 3{,}917 & 0.699 & 0.603 & $100\%$\\
48 & 9{,}808 & 13{,}895 & 15{,}475 & 0.706 & 0.634 & $100\%$\\
56 & 40{,}875 & 57{,}799 & 63{,}916 & 0.707 & 0.640 & $100\%$\\
60 & 77{,}805 & 110{,}839 & 127{,}157 & 0.702 & 0.612 & $100\%$\\
\bottomrule
\end{tabular}
\end{table}

比值在四个数量级的 $N$ 上保持平稳，正如常数因子理论所预测。同一内核上的 wall-clock 计时（表~\ref{tab:rho-wallclock}）确认该节省可端到端兑现：加入每步``昂贵预言机''负载 \textsc{extra}（模拟余因子工作或更重的混合映射）后，所提方法在每一档上都比 Brent 快约 $1.43\times$、比 Floyd 快约 $1.6\times$。

\begin{table}[H]\centering
\caption{每个分解出的 $48$ 位半素数的 wall-clock（$\mu$s，均值$\pm$标准差，每档 $400$ 个实例；Montgomery REDC，批量 gcd）。\textsc{extra} 为注入的每步预言机代价。}
\label{tab:rho-wallclock}
\begin{tabular}{rrrrcc}
\toprule
\textsc{extra} & Floyd & Brent & 所提 & 对 Brent & 对 Floyd\\
\midrule
0  & $150.2\pm84.2$ & $132.1\pm66.7$ & $92.3\pm45.4$ & $1.43\times$ & $1.63\times$\\
2  & $257.4\pm144$ & $226.6\pm113$ & $158.2\pm77$ & $1.43\times$ & $1.63\times$\\
8  & $629.8\pm500$ & $528.8\pm266$ & $369.7\pm181$ & $1.43\times$ & $1.70\times$\\
32 & $2064\pm1489$ & $1761\pm901$ & $1231\pm647$ & $1.43\times$ & $1.68\times$\\
64 & $3805\pm2132$ & $3360\pm1696$ & $2376\pm1213$ & $1.41\times$ & $1.60\times$\\
\bottomrule
\end{tabular}
\end{table}

\section{相关工作}
\label{sec:related}

Floyd 的龟兔算法通常归于其 1967 年关于非确定算法的工作~\cite{floyd1967};Brent 的二的幂方法
源于蒙特卡洛因子分解,据报道在 Pollard-$\rho$ 应用中相对 Floyd 减少约三分之一~\cite{brent1980},
我们的测量重现了这一点(第~\ref{sec:rho} 节)。Pollard 的 $\rho$~\cite{pollard1975} 是后继真正
昂贵的典型场景。随机映射的 $(\mu,\lambda)$ 分布由 Knuth--Flajolet 学派~\cite{flajolet,sedgewick2013}
刻画,其中 $\mathbb E[\mu+\lambda]\sim\sqrt{\pi/2}\sqrt n$。另一条路线以空间换求值:Gosper 的栈
方案~\cite{gosper1972} 与 Nivasch 的单趟栈~\cite{nivasch2004} 达到最优 $\mu+\lambda$,但使用
$\Theta(\log(\mu+\lambda))$ 内存;Sedgewick、Szymanski 与 Yao 的时空界~\cite{ssy1982} 表明该内存
是必要的。表~\ref{tab:related} 使用上文确立的闭式给出本方法的定位;其中
$k_B=\lceil\log_2\max(\lambda,\mu+1)\rceil$,$T=\lambda\lceil\mu/\lambda\rceil$。

\begin{table}[H]\centering
\caption{按辅助空间与后继求值次数划分的环检测方法。所提方法在保持 $O(1)$ 空间的同时逐点支配两个
常数空间基线;Nivasch 达到求值最优但使用非常数内存。}
\label{tab:related}
\begin{tabular}{llll}
\toprule
方法 & 辅助空间 & 后继求值次数(环形) & 备注\\
\midrule
Floyd~\cite{floyd1967} & $O(1)$ & $F=3T+2\mu$ & 龟/兔\\
Brent~\cite{brent1980} & $O(1)$ & $B=2^{k_B}\!-\!1+2\lambda+2\mu$ & 二的幂\\
\textbf{所提} & $O(1)$ & $\mu+\lambda\le A\le 2\mu+3\lambda$, $A\le\min(B,F)$ & 本文\\
Nivasch~\cite{nivasch2004} & $O(\log k)$ 期望 & $\mu+\lambda$(最优) & 栈\\
\midrule
下界(定理~\ref{thm:lowerbound}) & 任意 & $\ge\mu+\lambda$ & 本文\\
\bottomrule
\end{tabular}
\end{table}

\paragraph{昂贵后继与并行搜索。}
当后继是密码学映射时，另有三条相关路线。Pollard--Brent 对 $F_8$ 的分解~\cite{brent1981} 确立了 $\rho$ 方法的规模化，而 Montgomery 的批处理与求逆技巧~\cite{montgomery1987} 降低了我们保持固定的每步代价：它们与我们的调度正交且可组合，后者只改变后继被取用的\emph{次数}。Teske 的加法/混合游走~\cite{teske1998} 提升迭代映射的统计质量（缩小有效 $\mu+\lambda$），同样与检测调度正交。最后，Quisquater 与 Delescaille 的可区分点技术~\cite{qd1990} 以及 van Oorschot 与 Wiener 的并行碰撞搜索~\cite{vow1999} 以内存与多机换取吞吐；我们的贡献针对单次游走的\emph{串行}常数，而定位锚的可区分点变体是自然的延伸（第~\ref{sec:limits} 节）。

我们的贡献不是一个新的渐近类,而是一个更紧的常数,配以精确、可验证的分析、一个被证明的逐实例
支配性、一个匹配的下界,以及一个多锚点的延迟/工作量权衡。

\section{讨论与局限}
\label{sec:limits}
\begin{itemize}[leftmargin=2em,itemsep=2pt]
  \item 渐近复杂度仍为 $\Theta(\mu+\lambda)$;改进是常数因子(期望约 $2.18$ 对 $3.08$/$3.5$)。
  \item 支配结果只在环形输入上计后继求值。本方法比 Floyd 做更多指针比较,且比 Brent 控制流更
        复杂;在廉价后继情形,这些代价使它慢于朴素 Floyd(第~\ref{sec:wallclock} 节,
        表~\ref{tab:weighted})。
  \item 对无环链表,空检查是否计作一次解引用会改变精确常数,故我们只对环形输入主张解引用计数
        支配性,对其余情形仅主张 $O(n)$ 终止。
  \item 在某些应用(如 Pollard $\rho$)中,其他代价——批量 gcd、模约减——也可能重要,故端到端
        的节省须按应用重新测量。
  \item 本方法\emph{检测}复发,但其本身并不\emph{逃离}复发。在优化类场景中,它是一个廉价的
        振荡/极限环触发器,而非求解器。
  \item \textbf{代价模型。} 支配性度量只计后继求值。在比较与分支误预测并非免费的硬件上，净加速要求后继代价至少约为一次比较的 $8\times$（表~\ref{tab:weighted}，第~\ref{sec:wallclock} 节）；我们明确陈述这一门槛，而非主张无条件加速。
  \item \textbf{原创性。} 两阶段探测与二的幂锚点是 Brent 风格；据我们所知，新的是把一条\emph{失败压缩}（nextval）前缀跳过与一个\emph{定位锚点}（上一轮锚点，已证可安全地由此开始入口检测）相结合。正是这一跳过降低了\emph{总}求值次数（而不仅是检测），给出纯多锚方案可证无法达到的 $\approx2.18$ 常数（推论~\ref{cor:o1-optimal}）。
  \item \textbf{开放维度。} 可区分点/并行变体、平方的 SIMD 向量化批处理、在线/流式输入，以及对带噪声后继预言机的鲁棒性，在此均未触及；弥合已证 $O(1)$ 空间检测下界 $1.2213$ 与总常数 $2.18$ 之间的差距亦然。
\end{itemize}

\paragraph{定位。}
诚实的总结是：这是一次常数因子的精化，而非复杂度类的结果——$\Theta(\mu+\lambda)$ 时间与 $O(1)$ 空间不变。其价值集中于：(i) 后继主导代价之处（Pollard $\rho$、昂贵的混合映射），以及 (ii) 逐实例最坏情况保证重要之处，因为本方法在任一环形输入上都不输给 Brent 或 Floyd，且在几乎所有输入上获胜。

\section{结论}
\label{sec:conclusion}

我们提出了一种用于环检测的单调单指针锚点前进算法,它结合了两阶段探测、一条失败压缩(nextval)
跳过规则,以及一个从更靠近环处开始入口检测的定位锚点。我们证明了正确性、线性时间与常数空间;
给出了其后继求值次数的精确闭式;证明了它在\emph{每个}环形输入上都弱支配 Brent 与 Floyd——除三个
(分别为一个)退化输入外都严格支配——并且,对照一个匹配的 $\mu+\lambda$ 下界,它是常数因子最优的。
一项平均情形 Mellin 分析把它的常数确定为对数周期的 $\approx2.18033$,一个多锚点推广刻画了
检测/总工作量的 Pareto 前沿,而 wall-clock 与 Pollard-$\rho$ 研究定位出``更少求值''何时变为真实
加速的临界点。其结果是在后继求值模型下,对 Brent 风格环查找的一次干净、与分布无关、可完全分析的
精化——在后继是主导代价的场合最为有用——而非一个普适的替代品。

\begin{thebibliography}{9}

\bibitem{floyd1967} R.~W.~Floyd. Nondeterministic algorithms.
\emph{Journal of the ACM}, 14(4):636--644, 1967.

\bibitem{brent1980} R.~P.~Brent. An improved Monte Carlo factorization algorithm.
\emph{BIT Numerical Mathematics}, 20:176--184, 1980.

\bibitem{pollard1975} J.~M.~Pollard. A Monte Carlo method for factorization.
\emph{BIT Numerical Mathematics}, 15(3):331--334, 1975.

\bibitem{kmp1977} D.~E.~Knuth, J.~H.~Morris~Jr., and V.~R.~Pratt. Fast pattern
matching in strings. \emph{SIAM Journal on Computing}, 6(2):323--350, 1977.

\bibitem{flajolet} P.~Flajolet and A.~M.~Odlyzko. Random mapping statistics.
\emph{EUROCRYPT '89}, LNCS 434, 329--354, 1990.

\bibitem{nivasch2004} G.~Nivasch. Cycle detection using a stack.
\emph{Information Processing Letters}, 90(3):135--140, 2004.

\bibitem{gosper1972} R.~W.~Gosper. Item 132: cycle detection in $O(\log)$ space.
In M.~Beeler, R.~W.~Gosper, and R.~Schroeppel, \emph{HAKMEM}, MIT AI Memo 239, 1972.

\bibitem{ssy1982} R.~Sedgewick, T.~G.~Szymanski, and A.~C.~Yao. The complexity of
finding cycles in periodic functions. \emph{SIAM Journal on Computing},
11(2):376--390, 1982.

\bibitem{sedgewick2013} R.~Sedgewick and P.~Flajolet. \emph{An Introduction to
the Analysis of Algorithms}. Addison-Wesley, 2nd edition, 2013.

\bibitem{montgomery1987} P.~L.~Montgomery. Speeding the Pollard and elliptic
curve methods of factorization. \emph{Mathematics of Computation},
48(177):243--264, 1987.

\bibitem{brent1981} R.~P.~Brent and J.~M.~Pollard. Factorization of the eighth
Fermat number. \emph{Mathematics of Computation}, 36(154):627--630, 1981.

\bibitem{teske1998} E.~Teske. A space efficient algorithm for group structure
computation. \emph{Mathematics of Computation}, 67(224):1637--1663, 1998.

\bibitem{qd1990} J.-J.~Quisquater and J.-P.~Delescaille. How easy is collision
search? Application to DES. \emph{EUROCRYPT '89}, LNCS 434, 429--434, 1990.

\bibitem{vow1999} P.~C.~van Oorschot and M.~J.~Wiener. Parallel collision search
with cryptanalytic applications. \emph{Journal of Cryptology}, 12(1):1--28, 1999.

\end{thebibliography}

\end{document}
